\chapter{Diskussion}
\label{sec:diskussion}

Die in Kapitel 6 präsentierten Testergebnisse liefern ein klares Bild der Systemleistung. In diesem Kapitel werden diese Ergebnisse interpretiert und kritisch bewertet. Dabei wird untersucht, inwiefern die Forschungsziele erreicht wurde, wo die grenzen des Systems liegen und welche Funktionalitäten für die Qualität ausschlaggebend waren.

\section{Beantwortung der zentralen Forschungsfrage}
Die primäre Forschungsfrage dieser Arbeit war, ob eine Verarbeitungskette aus Objekterkennung, Bildsegmentierung, Objektverfolgung und zeitlicher Zustandsmodellierung, eine robustere Beute-Erkennung ermöglicht gegenüber Einzelbildmethoden.

Die Evaluierung liefert hierauf eine eindeutig Antwort. Ein einzelbildbasierter Ansatz ist bei der Analyse von Videostreams stark anfällig für Fehler. Bewegungsunschärfe (Motion Blur), wenn die Katze ins Bild läuft, partielle Verdeckungen (Okklusionen) der Beute durch die Pfoten oder rasche Änderungen der Beleuchtung führen zu starken Ungenauigkeiten der Segmentierung. Ohne die zeitliche Aggregation im Zustandsautomaten würde die Zugangskontrolle ständig zwischen den Zuständen \glqq verriegelt\grqq\ und \glqq entriegelt\grqq\ wechseln (Flackern).

Durch die hilfe der \texttt{StateMachine} in Kombination mit ByteTrack wird diese Problem elegant gelöst. Die zeitliche Aggregation über eine Historie von Einzelbildern erzeugt ein robustes Signal, das Ausreißer effektiv glättet. Die Ablationsstudie "Relevanz des Zustandsautomaten" zeigt, dass das System selbst dann richtig entscheidet, wenn einzelne Bilder falsch klassifiziert werden. Somit ist der Beweis erbracht, dass der Zustandsautomat in diesem Anwendungsfall ein wesentlicher Teil ist und zwingend für die Zuverlässigkeit im realen Betrieb ist.

\section{Bewertung des Architektur-Ansatzes}
Die Entscheidung, das Problem in zwei separate neuronale Netze (YOLO-Detektor für das Gesicht und YOLO-Segmentierungsmodell für den Zustand) zu unterteilen, hat sich als äußerst vorteilhaft erwiesen.

\subsection{Stärken der Verarbeitungsketten-Kaskadierung}
Der übliche Ansatz im Bereich maschinelles Sehen bestünde darin, ein einziges Objekterkennungsmodell darauf zu trainieren, direkt die Klassen \glqq Katze\grqq\ und \glqq Beutetier\grqq\ im Gesamtbild zu finden. Bei einer Katzenklappe scheitert dieser Ansatz jedoch häufig. Eine Maus im Maul einer Katze nimmt im Gesamtbild der Kamera nur einen verschwindend geringen Bruchteil der Pixel ein und ist durch das Maul und die Schnurrhaare stark verdeckt. Standard-Objektdetektoren haben immense Schwierigkeiten, derart kleine, okkludierte Objekte zuverlässig mit einem Begrenzungsrahmen zu versehen \cite{tong2020small}.

Der in dieser Arbeit gewählte Kaskadierungs-Ansatz umgeht dieses Problem elegant \cite{baumann2020smartcatflap, zhang2014partbased}. Der erste Detektor muss lediglich das vergleichsweise große und markante Katzengesicht finden – eine triviale Aufgabe für moderne Architekturen wie YOLO26, was sich in der hohen Erkennungsstabilität niederschlug. Durch das anschließende Herausschneiden (Zuschneiden) dieses Bereichs wird die Aufmerksamkeit des nachgelagerten Segmentierungsmodells exklusiv auf die relevante Region fokussiert (Aufmerksamkeitsmechanismus \textit{by design}). Das Segmentierungsmodell muss keine kleine Maus im Hintergrund des Gartens suchen, sondern analysiert lediglich die Pixelstruktur direkt am Maul der Katze.

\subsection{Die Rolle der asymmetrischen Randerweiterung}
Dadurch das die Annotationen des Datensatzes im Detektor (\glqq Cat Dataset\grqq\ von Crawford \cite{kaggle2024cat}) den unteren Teil des Mauls komplett verdecken, ist es unerlässlich für Training und Inferrenz den Begrenzungsrahmen nach untenhin Proportional zu erweitern.

\section{Schwächen, Limitationen und Fehlerquellen}
Aus der Analyse werden die Limitationen des aktuellen Systems genau aufgezeigt. Dadurch entsteht das Potenzial für zukünftige Verbesserungen:

\begin{enumerate}
    \item \textbf{Verlust von Tiefeninformationen:} Dadurch das die Kamera ein 2D Bild aufnimmt kann nicht unterschieden werden, wie weit weg sich ein Objekt unterhalb des Katzenmauls befindet. Daraus folgt, dass das Modell den Unterschied zwischen Hintergrund und Beute nicht zuverlässig trennen kann.
    \item \textbf{Extreme Bewegungsgeschwindigkeit:} Wenn eine Katze fluchtartig durch die Klappe rennt, bleibt sehr wenig Zeit eine klare Aussage zu treffen. Bei einer Laufzeit von ca. 8 FPS verbleiben dem System nur eine Handvoll Einzelbilder zur Analyse.  In diesem Szenario kann die \texttt{StateMachine} ihre Stärke (die zeitliche Glättung) nicht voll ausspielen und die Katze muss warten bis der Zustandsautomat eine klare Entscheidung trifft.
    \item \textbf{Farbe der Katze:} Wenn eine Katze ein sehr dunkles Fell besitzt, hat die Segmentierung Schwierigkeiten und verliert an Genauigkeit. Hier ist es essentiell Daten aus dem Praxisbetrieb zu sammeln, bei denen sich das Modell unsicher war, um genau diese Fälle besser trainieren zu können.
\end{enumerate}

\section{Lokale Datenverarbeitung}
Ziel dieser Arbeit war es auch, dass das System ressourcenschonend konzipiert wird um sicherzustellen dass die Inferrenz auch auf einem Einplatinencomputer lokal operieren kann. Daher viel die Entscheidung auf eine leichtgewichtige Architektur (YOLO26n) und optimierte Formate (TensorFlow Lite). Die dadurch verlorene Genauigkeit wird durch die erreichte Leistung gerechtfertigt (ca. 8 FPS).

Dieser Ansatz beweist, dass eine komplexe Aufgabe auf ressourcenbegrenzter Hardware durch eine intelligent gewählte Architektur gelöst werden kann, ohne dass die Robustheit zu stark beeinträchtigt wird. Dadurch wird eine Auslagerung von Berechnungen in externe Rechenzentren vermieden. Das spart Geld und macht die Lösung zuverlässiger, da keine Internetverbindung notwendig ist.

Abschließend lässt sich sagen, dass die Architektur ihre Aufgabe größtenteils erfüllt. Die verbleibenden Fehlalarme sind kein Architekturfehler, sondern zeigen die Grenzen des begrenzten Segmentierungsdatensatzes auf. Mithilfe des in Kapitel 5 integrierten Modells zum Speichern von Bildern, bei denen das Modell unsicher war, können genau diese Lücken im Datensatz ergänzt werden. So kann das Modell kontinuierlich verbessert werden.